\section{Trees and more graphs}

\subsection{Tree facts}\label{t:trees}
\begin{tip}
\begin{itemize}
\item $n$ nodes, $n - 1$ edges. Connected with $n - 1$ edges $\Rightarrow$ a tree.
\item Exactly one path between two nodes: one BFS gives all distances, also weighted.
\item Connected graph: $m - (n - 1)$ extra edges, each closes a cycle. Forest of $c$ trees:
  $n - c$ edges. Few extra edges: try all $2^k$ ways to handle them.
\item Edge above $v$ is used by $\mathit{sz}[v] \cdot (n - \mathit{sz}[v])$ paths. Summed over
  all edges: the sum of all pair distances.
\item Parent list with $p_i < i$: loop $i = n, \dots, 2$ for children before parents.
\item Trees are bipartite: colour by depth mod 2.
\end{itemize}
\end{tip}

\topic{Rooting a tree}{$O(n)$}
\use{Parent, depth, \texttt{order}. Backwards over \texttt{order}: children before parents
(sizes, tree DP). Forwards: parents first (rerooting). No recursion.}
\code{tree_root}
\code{tree_size}

\topic{DFS order}{$O(n)$}
\use{Enter / leave times. $u$ is in the subtree of $v$ iff
\texttt{tin[v] <= tin[u] < tout[v]}: a subtree is a range of an array.}
\code{tree_dfs}

\topic{Tree DP}{$O(n)$}
\use{The answer for $v$ follows from its children's answers. Fill backwards over
\texttt{order}; the answer is at the root.}
\code{tree_dp}

\topic{Diameter}{$O(n)$}
\use{Longest path: BFS from any node to the farthest $a$, BFS from $a$ to the farthest $b$.
The farthest node from any $v$ is $a$ or $b$.}
\code{tree_diameter}

\topic{LCA and binary lifting}{$O(\log n)$ per query}\label{t:lca}
\use{Distance between $a$ and $b$, the $k$-th node above $v$. Weighted:
$w[a] + w[b] - 2\,w[\mathrm{lca}]$. Every node has one next node (functional graph): where
after $k \le 10^{18}$ steps? \texttt{Jump}.}
\code{tree_lca}
\code{graph2_jump}

\topic{Rerooting}{$O(n)$}
\use{A tree DP answer for every node as the root. Moving the root from \texttt{par[v]} to
$v$: $\mathit{sz}[v]$ nodes get 1 closer, the other $n - \mathit{sz}[v]$ 1 farther.}
\code{tree_reroot}

\topic{Strongly connected components}{$O(n + m)$}\label{t:scc}
\use{Directed graph, nodes that reach each other become one node: the condensation is a DAG.
To reach everything you need a start in every component with in-degree 0.}
\code{graph2_scc}

\topic{Euler path}{$O(n + m)$}
\use{Every edge exactly once. Undirected: 0 or 2 odd-degree nodes (start at one).
Directed: in $=$ out everywhere, or one start (out $=$ in $+ 1$) and one end. Circuit: all
even / in $=$ out.}
\code{graph2_euler}

\topic{Bipartite check}{$O(n + m)$}
\use{Two groups, every edge between them: possible iff no odd cycle.}
\code{graph2_bipartite}

\topic{Bipartite matching}{$O(n \cdot m)$}
\use{Most pairs between two kinds of things. K\"onig: min vertex cover $=$ max matching;
max independent set $= n -$ max matching.}
\code{graph2_matching}
